\documentclass[a4paper,12pt]{article}

\usepackage[T2A]{fontenc}
\usepackage[utf8]{inputenc}
\usepackage[russian]{babel}
\usepackage{amsmath,amssymb,mathtools}
\usepackage{helvet}
\renewcommand{\familydefault}{\sfdefault}
\usepackage{icomma}
\usepackage{geometry}
\usepackage{microtype}
\usepackage{tikz}
\usetikzlibrary{calc,arrows.meta,positioning,shapes.geometric}
\usepackage{pgfplots}
\pgfplotsset{compat=1.18}
\usepackage{tabularx,booktabs}
\usepackage{enumitem}
\usepackage{hyperref}
\usepackage{parskip}
\usepackage{fancyhdr}
\usepackage{setspace}
\usepackage{xcolor}

\geometry{margin=2cm}
\onehalfspacing

\definecolor{accent}{HTML}{006666}
\definecolor{darkaccent}{HTML}{003F3F}
\definecolor{textgray}{HTML}{333333}

\pagestyle{fancy}
\fancyhf{}
\fancyhead[L]{Физика. Кинематика}
\fancyhead[R]{Иван Петраченков}
\fancyfoot[C]{\thepage}

\begin{document}

\begin{titlepage}
\centering

{\Large\textcolor{darkaccent}{Чек-лист}\par}

\vspace{1cm}

{\huge Кинематика:\\
путь, перемещение и описание движения\par}

\vspace{1cm}

{\Large ЕГЭ. Путь А\par}

\vfill

{\large Лёвин Артём Александрович\\
эксклюзивно для Ивана Петраченкова\par}

\vspace{0.5cm}

{\large \today}

\vspace{1cm}

\begin{tikzpicture}
\draw[accent,line width=1pt]
(0,0)
.. controls (2,0.8) and (5,-0.8) ..
(9,0);
\end{tikzpicture}

\end{titlepage}

\section*{1. Базовые понятия кинематики}

\textbf{Материальная точка} --- модель тела, размерами которого в данной задаче можно пренебречь.

\textbf{Траектория} --- линия, по которой движется тело.

\textbf{Путь $S$} --- длина траектории.

\textbf{Перемещение $\vec{s}$} --- вектор, соединяющий начальное и конечное положения тела.

Главное различие:

\[
S\geq |\vec{s}|
\]

Путь всегда неотрицателен, а перемещение является вектором.

\subsection*{Частные случаи}

Если тело вернулось в исходную точку:

\[
|\vec{s}|=0
\]

но путь может быть отличен от нуля.

Например, полный оборот по окружности радиуса $R$:

\[
S=2\pi R,\qquad |\vec{s}|=0
\]

Если тело переместилось в противоположную точку окружности:

\[
S=\pi R,\qquad |\vec{s}|=2R
\]

\section*{2. Координаты и векторы}

Для вектора:

\[
\vec{AB}(x_B-x_A;\ y_B-y_A)
\]

Длина вектора:

\[
|\vec{AB}|=
\sqrt{(x_B-x_A)^2+(y_B-y_A)^2}
\]

При движении вдоль оси $Ox$:

\[
s_x=x-x_0
\]

где $x_0$ --- начальная координата.

\subsection*{Проекции вектора}

Если вектор длины $A$ образует угол $\alpha$ с осью $Ox$:

\[
A_x=A\cos\alpha
\]

\[
A_y=A\sin\alpha
\]

Запомнить:

\[
\cos\alpha=
\frac{\text{прилежащий катет}}
{\text{гипотенуза}}
\]

\[
\sin\alpha=
\frac{\text{противолежащий катет}}
{\text{гипотенуза}}
\]

\section*{3. Равномерное движение}

При равномерном движении скорость постоянна.

Основное уравнение:

\[
x=x_0+v_xt
\]

где:

\begin{itemize}
\item $x_0$ --- начальная координата;
\item $v_x$ --- проекция скорости;
\item $t$ --- время.
\end{itemize}

Связь скорости и координаты:

\[
v_x=\frac{\Delta x}{\Delta t}
\]

Знак скорости показывает направление:

\[
v_x>0
\]

--- движение в положительном направлении оси.

\[
v_x<0
\]

--- движение против направления оси.

\section*{4. График координаты}

Формула

\[
x=x_0+v_xt
\]

имеет вид линейной функции.

На графике:

\begin{itemize}
\item горизонтальная ось --- время $t$;
\item вертикальная ось --- координата $x$.
\end{itemize}

Начальная координата:

\[
x_0=x(0)
\]

Скорость определяется наклоном графика:

\[
v_x=\tg\alpha
\]

где $\alpha$ --- угол наклона прямой к оси времени.

\section*{5. Средняя скорость}

Средняя скорость:

\[
v_{\text{ср}}
=
\frac{S_{\text{общ}}}{t_{\text{общ}}}
\]

Нельзя просто считать:

\[
\frac{v_1+v_2}{2}
\]

если времена движения различаются.

\subsection*{Пример}

Половину времени тело движется со скоростью $10$ км/ч, половину времени --- со скоростью $30$ км/ч.

Тогда:

\[
v_{\text{ср}}
=
\frac{10+30}{2}
=
20
\text{ км/ч}
\]

Но если половины относятся к пути, формула изменяется.

\section*{6. Протяжённые тела}

Если объект нельзя считать материальной точкой:

\[
t=\frac{L_1+L_2}{v_1+v_2}
\]

при движении навстречу.

При движении вдогонку:

\[
t=\frac{L_1+L_2}{|v_1-v_2|}
\]

Важно:

\[
L_1,\ L_2
\]

--- длины объектов.

Это не пути, которые они проходят.

\clearpage

\section*{Типовой алгоритм решения задач на путь и перемещение}

\begin{tikzpicture}[
node distance=16mm,
>=Latex,
every node/.style={font=\small,align=center},
start/.style={ellipse,draw=accent,minimum width=3cm},
process/.style={rounded rectangle,draw=accent,minimum width=4cm},
decision/.style={diamond,draw=accent,aspect=2},
arrow/.style={->,thick}
]

\node[start](s){Начало};

\node[process,below=of s](a)
{Определить\\начальное и конечное положение};

\node[decision,below=of a](b)
{Начало и конец\\совпадают?};

\node[process,below left=18mm and 15mm of b](c)
{Перемещение\\равно нулю};

\node[process,below right=18mm and 15mm of b](d)
{Построить вектор\\из начала в конец};

\node[process,below=30mm of b](e)
{Найти длину\\перемещения};

\node[start,below=of e](f)
{Ответ};

\draw[arrow](s)--(a);
\draw[arrow](a)--(b);
\draw[arrow](b)--node[left]{да}(c);
\draw[arrow](b)--node[right]{нет}(d);
\draw[arrow](c)--(e);
\draw[arrow](d)--(e);
\draw[arrow](e)--(f);

\end{tikzpicture}

\clearpage

\section*{Тренировочный блок}

\begin{enumerate}[itemsep=8pt]

\item Тело совершило полный оборот по окружности радиуса $5$ м.
Найдите путь и модуль перемещения.

\item Автомобиль проехал половину окружности радиуса $20$ м.
Найдите путь и модуль перемещения.

\item Точка переместилась из $(2;3)$ в $(8;11)$.
Найдите координаты перемещения.

\item Найдите длину вектора $(6;8)$.

\item Вектор длиной $12$ м образует угол $60^\circ$ с осью $Ox$.
Найдите его проекцию на ось $Ox$.

\item Координата тела меняется по закону:

\[
x=4+5t
\]

Найдите начальную координату и скорость.

\item За $20$ секунд координата тела изменилась на $60$ м.
Найдите скорость.

\item Прямая на графике $x(t)$ образует угол $30^\circ$ с осью времени.
Найдите скорость, если единичный масштаб выбран так, что
$\tg30^\circ=\frac1{\sqrt3}$.

\item Половину времени автомобиль двигался со скоростью
$40$ км/ч, вторую половину времени --- со скоростью
$80$ км/ч.
Найдите среднюю скорость.

\item Первую половину пути автомобиль прошёл со скоростью
$30$ км/ч, вторую --- со скоростью $60$ км/ч.
Найдите среднюю скорость.

\item Поезд длиной $400$ м проходит мост длиной $200$ м
со скоростью $72$ км/ч.
Найдите время движения.

\item Два тела движутся навстречу друг другу со скоростями
$6$ м/с и $4$ м/с.
Расстояние между ними $300$ м.
Найдите время встречи.

\item Тело движется по закону:

\[
x=-10+2t
\]

Через сколько секунд оно окажется в начале координат?

\item Человек обошёл озеро диаметром $100$ м.
Найдите путь и модуль перемещения.

\end{enumerate}

\clearpage

\section*{Ответы}

\renewcommand{\arraystretch}{1.3}

\begin{center}
\begin{tabularx}{\linewidth}{cX}
\toprule
№ & Ответ\\
\midrule
1 & $S=10\pi$ м, $|\vec{s}|=0$\\
2 & $S=20\pi$ м, $|\vec{s}|=40$ м\\
3 & $(6;8)$\\
4 & $10$\\
5 & $6$ м\\
6 & $x_0=4,\ v=5$ м/с\\
7 & $3$ м/с\\
8 & $\frac1{\sqrt3}$ м/с\\
9 & $60$ км/ч\\
10 & $40$ км/ч\\
11 & $30$ с\\
12 & $30$ с\\
13 & $5$ с\\
14 & $100\pi$ м, $0$\\
\bottomrule
\end{tabularx}
\end{center}

\clearpage

\end{document}